\section{Generalized Hodge and arithmetic consequences}
\label{cc:section}

Theorem~\ref{thm:cm} applies to products of CM abelian varieties and
therefore algebraizes rational Hodge correspondences between them.
We use this to obtain the algebraic supports predicted by the
generalized Hodge conjecture, following Abdulali's effective-realization
argument
\cite[Lemma~1 and Theorems~3--4]{Abdulali2005CM}. Hazama had earlier
reduced this implication to codimension-two Hodge classes on a special
family of CM abelian varieties
\cite[Theorem~8.3]{Hazama2003CM}. We then apply Milne's Hodge-to-Tate
and polarization theorems to obtain the Tate and Hodge standard
conjectures for abelian varieties.
A final application uses CM transport to algebraize specialized Hodge
classes on arbitrary good-reduction abelian varieties.

\subsection{The generalized Hodge conjecture}

A pure rational Hodge structure is \emph{effective} if all its nonzero
Hodge components have nonnegative indices. Recall that
\(V(r)^{p,q}=V^{p+r,q+r}\), so the twist \(V(r)\) has weight
\(\operatorname{wt}(V)-2r\). The Hodge circle of a pure Hodge structure
\(V\) acts on \(V^{p,q}\) by \(z^{p-q}\), for \(|z|=1\). Let \(J(V)\)
be the smallest algebraic subgroup of \(\GL(V)\) over \(\Q\) whose real
points contain that circle. A polarizable \(V\) is \emph{of CM type}
if \(J(V)\) is a torus, with the trivial group allowed. Every Hodge
substructure of the cohomology of a CM abelian variety is of CM type:
its Hodge-circle group is an image of the torus acting on degree-one
cohomology.

For a smooth projective complex variety \(X\), the geometric coniveau
filtration is
\[
 N^rH^k(X,\Q)=
 \sum_{\substack{Z\subset X\text{ closed}\\ \codim_X Z\ge r}}
 \ker\bigl(H^k(X,\Q)\longrightarrow H^k(X\setminus Z,\Q)\bigr).
\]
Grothendieck's corrected generalized Hodge conjecture says that a
rational Hodge substructure \(W\subset H^k(X,\Q)\) lies in
\(N^rH^k(X,\Q)\) whenever \(W(r)\) is effective; see
\cite[\S2]{Grothendieck1969General} and \cite[\S2]{Hazama2003CM}.
Equivalently, one closed subset of
codimension at least \(r\) supports all of \(W\). Indeed, finite
dimensionality permits taking a finite union of the supports needed
for a basis of \(W\).

We will obtain such a support by algebraizing a Hodge map from an
auxiliary CM abelian variety.  To get the required codimension, that
variety must realize an effective twist of \(W\) in cohomology of the
same weight; the twisting index will determine the codimension of the
resulting correspondence.  The next lemma constructs it by expressing
the Hodge indices at each embedding as sums of weight-one indices.

\begin{lemma}[Effective realization of a simple CM Hodge structure]
\label{cc:realization}
Let \(W\) be a nonzero simple polarizable rational Hodge structure of
CM type and weight \(k\). Put
\[
 r_0=\min\{p:W^{p,k-p}\ne0\},\qquad n=k-2r_0.
\]
Then \(n\ge0\), and there is a CM abelian variety \(B\), possibly a
point, and a surjective rational Hodge map
\[
 H^n(B,\Q)(-r_0)\longrightarrow W.
\]
\end{lemma}

\begin{proof}
Set \(U=W(r_0)\). Hodge symmetry shows that \(U\) is effective of
weight \(n\ge0\), with smallest Hodge index zero. If \(n=0\), every
type of \(U\) is \((0,0)\). Simplicity then gives \(U\simeq\Q\), and
we take \(B\) to be a point.

Suppose \(n>0\).  We first show that \(U\) is a line over a CM field;
we will then decompose its Hodge type into weight-one CM types.
Write \(T=J(U)\). Let \(E_0\) be the
\(\Q\)-subalgebra of \(\End_\Q(U)\) generated by the image of \(T\).
After splitting the torus, its action is diagonal by characters, so
\(E_0\) is a commutative finite étale algebra. Its field idempotents
and every subspace linear over one field factor are \(T\)-stable.
Simplicity of \(U\) therefore forces \(E_0\) to be a field \(E\) and
\(\dim_E U=1\).

We check that \(E\) is a CM field. Let the Weil operator \(C_U\) act
as \(i^{p-q}\) on \(U^{p,q}\). By polarizability there is a rational
Hodge pairing \(\psi\), with the conventional weight-dependent sign
absorbed in its definition, for which
\[
 q(x,y)=\psi(C_Ux,y)
\]
is positive definite on \(U_\R\). The adjoint \(\dagger\) for
\(\psi\) sends \(t\in T\) to \(t^{-1}\), hence preserves \(E\).
Every element of \(E\) commutes with \(C_U\), so \(\dagger\) is also
the adjoint for \(q\). On the real finite étale algebra \(E_\R\), it
must fix each primitive idempotent: if an idempotent \(e\) were sent
to a different one, then \(e^\dagger e=0\), contrary to
\(q(ex,ex)=q(x,e^\dagger ex)>0\) for some \(x\). On a complex factor
it must act by complex conjugation, since the identity would give
\(i^\dagger i=-1\), again contrary to positivity. Thus if
\(\dagger\) is trivial on \(E\), the field is totally real. If it is
nontrivial, a real embedding would force \(\dagger\) to be trivial
on \(E\), so there are no real embeddings; the fixed algebra in
\(E_\R\) is then a product of copies of \(\R\). Its fixed field is
therefore totally real, and \(E\) is CM.

The totally real case cannot occur here. The Hodge circle acts
through \(E_\R^\times\); its image in each real scalar factor is a
compact connected subgroup of \(\R^\times\), hence is trivial.
All Hodge types of \(U\) would then have equal indices. Since their
smallest index is zero, this would force \(n=0\). Consequently \(E\)
is CM, and \(\dagger\) is its complex conjugation.

Let \(\Sigma=\Hom(E,\C)\), and write \(\bar\sigma\) for the conjugate
of \(\sigma\). Since \(E\) acts by Hodge endomorphisms and \(U\) is an
\(E\)-line,
\[
 U_\C=\bigoplus_{\sigma\in\Sigma}U_\sigma,\qquad
 U_\sigma\text{ has type }(a_\sigma,n-a_\sigma),
\]
where \(0\le a_\sigma\le n\) and
\(a_\sigma+a_{\bar\sigma}=n\). Choose \(n\) CM types
\(\Phi_1,\ldots,\Phi_n\) of \(E\) with
\[
 \#\{j:\sigma\in\Phi_j\}=a_\sigma.
\]
This is possible independently for each conjugate pair: put
\(\sigma\) into any \(a_\sigma\) lists and \(\bar\sigma\) into the
remaining lists.

For each \(j\), give the rational vector space \(H_j=E\) the
weight-one Hodge structure whose \(\sigma\)-line has type \((1,0)\)
exactly when \(\sigma\in\Phi_j\).  The polarization and Riemann
construction of Section~\ref{sec:tensors}, surrounding
Equation~\eqref{eq:cm-types}, applies to any CM field, with no
Galois hypothesis.  It gives a CM abelian variety \(B_j\) with
\(H^1(B_j,\Q)\simeq H_j\).

Identify \(U\) and the \(H_j\) with \(E\) as \(E\)-lines. The rational
multiplication map
\[
 E^{\otimes_\Q n}\longrightarrow E,\qquad
 x_1\otimes\cdots\otimes x_n\longmapsto x_1\cdots x_n
\]
is surjective. Over \(\C\) it kills tensors with unequal embedding
labels, while the tensor with all labels equal to \(\sigma\) has type
\((a_\sigma,n-a_\sigma)\) and maps to \(U_\sigma\). It is therefore a
Hodge map to \(U\). For \(B=\prod_{j=1}^nB_j\), its source is the
rational Künneth summand \(\bigotimes_j H^1(B_j,\Q)\) of \(H^n(B,\Q)\).
Extending by zero on the other summands and twisting by \((-r_0)\)
proves the lemma.
\end{proof}

\begin{corollary}[Generalized Hodge conjecture for CM abelian varieties]
\label{cc:ghc}
Let \(A\) be a complex CM abelian variety. For integers \(k,r\ge0\)
and a rational Hodge substructure \(W\subset H^k(A,\Q)\) such that
\(W(r)\) is effective, there is a closed algebraic subset \(Z\subset A\)
of codimension at least \(r\) such that
\[
 W\subset\ker\bigl(H^k(A,\Q)\longrightarrow
                         H^k(A\setminus Z,\Q)\bigr).
\]
This includes every finite product of CM abelian varieties and every
power of such a product.
\end{corollary}

\begin{proof}
The zero substructure permits \(Z=\varnothing\). Otherwise,
semisimplicity of polarizable rational Hodge structures reduces the
claim to a simple summand \(W\), after which we take the finite union
of the resulting supports. Let \(r_0\) be its smallest Hodge index.
Then \(r_0\ge r\). Also \(r_0\le k/2\le\dim A\), since \(W\ne0\).
Put \(n=k-2r_0\), take \(B\) from
Lemma~\ref{cc:realization}, and write \(b=\dim B\). Composing the map
of that lemma with \(W\hookrightarrow H^k(A,\Q)\) gives
\[
 f:H^n(B,\Q)(-r_0)\longrightarrow H^k(A,\Q),
 \qquad \im(f)=W.
\]
By Poincaré duality, \(f\) corresponds to a rational Hodge class of
type \((b+r_0,b+r_0)\) in
\[
 H^{2b-n}(B,\Q)\otimes H^k(A,\Q)
 \subset H^{2b+2r_0}(B\times A,\Q).
\]
The product \(B\times A\) is CM. Theorem~\ref{thm:cm} therefore
represents this exact class by
\(\Gamma\in\CH^{b+r_0}(B\times A)_\Q\), whose action in degree \(n\)
is \(f\).

Choose a finite rational cycle representing \(\Gamma\), and let
\(Z\) be the projection of its support to \(A\). Every component of
that cycle has dimension
\[
 \dim(B\times A)-(b+r_0)=\dim A-r_0.
\]
The projection is proper, so \(Z\) is closed and
\(\codim_A Z\ge r_0\ge r\). The correspondence restricts to zero on
\(B\times(A\setminus Z)\); functoriality of its cohomological action
therefore makes its image \(W\) restrict to zero on \(A\setminus Z\).
\end{proof}

\subsection{The Tate conjecture over finite fields}

Let \(A\) be an abelian variety over \(\Fq\), put
\(\bar A=A\times_{\Fq}\overline{\Fq}\), and write
\(G_q=\Gal(\overline{\Fq}/\Fq)\). For a prime
\(\ell\ne\operatorname{char}(\Fq)\), the rational Tate conjecture in
codimension \(r\) asks whether cycles over \(\Fq\) span
\(H^{2r}_{\mathrm{\acute et}}(\bar A,\Q_\ell(r))^{G_q}\).

\begin{corollary}[Tate conjecture for abelian varieties over finite fields]
\label{cc:tate}
For every abelian variety \(A/\Fq\), every integer \(r\ge0\), and every
prime \(\ell\ne\operatorname{char}(\Fq)\), the cycle-class map
\[
 \CH^r(A)\otimes_\Z\Q_\ell\longrightarrow
 H^{2r}_{\mathrm{\acute et}}(\bar A,\Q_\ell(r))^{G_q}
\]
is surjective. Numerical equivalence and \(\ell\)-adic homological
equivalence coincide on rational algebraic cycles on \(\bar A\).
\end{corollary}

\begin{proof}
Apply Milne's Theorem~7.1 and his formulation (0.1)
\cite{Milne1999Tate}, with the correction in
\cite{cc:Milne1999Addenda}: the cofinal auxiliary Galois CM fields
are enlarged to contain a real subfield of degree \(>2\) inert at
the residue prime. The theorem's hypothesis is supplied by
Theorem~\ref{thm:cm}. It gives the stated equality of
equivalence relations and algebraic spanning of the geometric Tate
classes: those fixed over some finite field of definition.

Let \(\alpha\) lie in the displayed \(G_q\)-invariant space. The
geometric conclusion writes \(\alpha\) as a finite
\(\Q_\ell\)-linear combination of cycle classes. All the cycles in
this expression descend to \(A_{\mathbb F_{q^m}}\) for one \(m\);
let \(z_m\in\CH^r(A_{\mathbb F_{q^m}})\otimes_\Z\Q_\ell\) denote
the resulting combination. For the finite projection
\(\pi_m:A_{\mathbb F_{q^m}}\to A\), proper pushforward and base
change for cycle classes give
\[
 \operatorname{cl}_\ell\!\left(\frac1m(\pi_m)_*z_m\right)
 =\frac1m\sum_{\sigma\in
        \Gal(\mathbb F_{q^m}/\Fq)}\sigma(\alpha)=\alpha.
\]
Thus \(m^{-1}(\pi_m)_*z_m\) is the required class over \(\Fq\).
Division by \(m\) is valid in \(\Q_\ell\), even when \(m\) is
divisible by \(\ell\) or by the characteristic.
\end{proof}

Milne's implication compares CM and finite-field abelian motives. It
does not assert that each Tate class lifts to a Hodge class on a
single CM lifting, even up to isogeny
\cite[Introduction, Remark~(a)]{Milne1999Tate}.

\subsection{The Hodge standard conjecture}

Let \(A\) be an abelian variety of dimension \(g\) over an
algebraically closed field \(K\), and let
\(A^r_{\mathrm{num}}(A)_\Q\) be the space of rational
codimension-\(r\) cycles modulo numerical equivalence. A
\emph{rational ample divisor class} is an
\(h\in A^1_{\mathrm{num}}(A)_\Q\) for which some positive integral
multiple is the numerical class of an ample Cartier divisor. For
integers \(0\le r\le g/2\), put
\[
 P_h^r(A)=\ker\!\left(
 h^{g-2r+1}:A^r_{\mathrm{num}}(A)_\Q
 \longrightarrow A^{g-r+1}_{\mathrm{num}}(A)_\Q\right).
\]

\begin{corollary}[Hodge standard conjecture for abelian varieties]
\label{cc:standard}
Let \(K\) be an algebraically closed field and \(A/K\) an abelian
variety of dimension \(g\). For every prime
\(\ell\ne\operatorname{char}(K)\), numerical
and \(\ell\)-adic homological equivalence coincide on rational
algebraic cycles on \(A\). For every rational ample divisor class
\(h\) and every integer \(0\le r\le g/2\), the rational symmetric form
\[
 P_h^r(A)\times P_h^r(A)\longrightarrow\Q,\qquad
 (x,y)\longmapsto(-1)^r\deg(xy\,h^{g-2r})
\]
is positive definite after extension of scalars to \(\R\). Thus the
Hodge standard conjecture for \(\ell\)-adic étale cohomology holds
for abelian varieties in arbitrary characteristic.
\end{corollary}

\begin{proof}
Apply Milne's Theorem~3.3 \cite{Milne2002Polarizations} using
Theorem~\ref{thm:cm}. In Milne's proof, the same Hodge hypothesis
supplies the compatibility and positivity of the relevant
polarizations through Theorem~2.1, read with the author's addendum
\cite{cc:Milne2002Addendum} and the corrected auxiliary-field choice
from \cite{cc:Milne1999Addenda}. Theorem~3.3 gives both the equality
of equivalence relations and positivity on the primitive space of
cycles modulo homological equivalence. The former identifies that
primitive space with the numerical one above.

The case \(g=0\) is immediate. For \(g>0\), choose a positive integer
\(m\) so that \(mh\) is the class of a smooth hyperplane section;
clearing denominators, taking a sufficiently large very ample
multiple, and Bertini's theorem permit this. Multiplication by the
nonzero rational scalar \(m^{g-2r+1}\) does not change the primitive
kernel, and the form for \(mh\) is \(m^{g-2r}\) times the form for
\(h\). This positive factor preserves positive definiteness.
\end{proof}

\subsection{Specialization of Hodge classes}

\begin{remark}[Algebraicity of specialized Hodge classes]
\label{cc:algebraic-specialization}
Theorem~\ref{thm:cm} also combines with the CM transport of
\cite[Corollary~3.4]{OpenAIMilneRationality2026} to give a
specialization statement for abelian varieties that need not be CM.
Fix a prime \(p\) and a \(p\)-adic place of \(\Qbar\).
If \(A/\Qbar\) has good reduction \(A_0/\overline{\mathbb F}_p\)
at that place, every rational Hodge class on \(A\) specializes to
the class of one rational algebraic cycle on \(A_0\), simultaneously
in all prime-to-\(p\) realizations and in crystalline cohomology
after extension to the completed field at the chosen place.
Indeed, the transport writes that specialization as the pullback of
a specialized CM Hodge class by a rational homomorphism of special
fibers. Theorem~\ref{thm:cm} algebraizes the CM class; specialization
and Chow pullback after clearing the homomorphism's denominator give
the required cycle. See
\cite[Corollary~6.4]{OpenAIMilneRationality2026} for the
field-of-definition argument and precise realization conventions,
including \(p=2\). This application uses the CM transport, not the
independent rationality proof of that paper, and is not an input to
Theorem~\ref{thm:cm}. It asserts algebraicity on the special fiber,
not the Hodge conjecture for arbitrary complex abelian varieties.
\end{remark}
